% ======================= CHAPTER 5: CONCLUSION ==========================
\chapter{Conclusion and Future Work}

\section{Summary of Contributions}

This thesis presented a systematic comparative evaluation of four deep learning architectures for multi-class skin lesion classification on the ISIC 2018 dataset. The key contributions of this work are:

\begin{enumerate}
    \item \textbf{An optimized training pipeline} that effectively addresses the severe class imbalance (58:1 ratio) inherent in dermatological datasets through a combination of Weighted Random Sampling, Mixup, CutMix, and Label Smoothing. This pipeline demonstrated consistent effectiveness across all four architectures, eliminating overfitting in every case.
    
    \item \textbf{A comprehensive architectural comparison} spanning three CNN families (EfficientNet-B0, ResNet50, DenseNet121) and one Vision Transformer (Swin-Tiny). The results revealed that EfficientNet-B0 achieves the best individual performance (88.36\% accuracy, 0.831 macro F1) with the fewest parameters (5.3M), while Swin-Tiny provides the best probability calibration (0.960 AUC) but suffers from limited data availability.
    
    \item \textbf{A soft-voting ensemble} that leverages architectural diversity to boost performance to 89.49\% accuracy and 0.980 macro AUC, surpassing all individual models and demonstrating the complementary nature of CNNs and Transformers.
    
    \item \textbf{Grad-CAM interpretability analysis} confirming that all models attend to clinically relevant lesion features, a critical requirement for any AI system intended for medical deployment.
    
    \item \textbf{An ablation study} quantifying the individual contribution of each regularization technique, providing empirical evidence for the design decisions underlying the pipeline.
\end{enumerate}

\section{Key Findings}

The experimental results yielded several insights with practical implications:

\textbf{Efficiency matters:} EfficientNet-B0, despite being 5$\times$ smaller than ResNet50 and Swin-Tiny, consistently outperformed them on hard-classification metrics. This challenges the notion that larger models necessarily produce better results, particularly in data-constrained medical imaging scenarios.

\textbf{CNNs vs. Transformers:} An interesting inverse relationship was observed between accuracy and AUC across the four models. CNNs (particularly EfficientNet) excelled at discrete classification, while the Transformer (Swin-Tiny) produced better-calibrated probabilities. This suggests that in clinical applications where probability ranking is more important than binary decisions (e.g., triage systems), Transformers may be preferred despite lower accuracy.

\textbf{Melanoma remains the hardest class:} Across all models and the ensemble, melanoma (MEL) consistently achieved the lowest F1-scores, reflecting the genuine clinical difficulty of distinguishing early-stage melanoma from benign nevi. This highlights the need for continued research into melanoma-specific feature learning.

\textbf{Ensembles are effective:} The simple averaging of probability distributions from four models provided a meaningful improvement across all metrics, with the most notable gain in AUC (+2.8\%). This improvement comes at no additional training cost, requiring only inference from multiple models.

\section{Limitations}

This study has several limitations that should be acknowledged:

\begin{itemize}
    \item \textbf{Single dataset:} All experiments were conducted on ISIC 2018. Cross-dataset evaluation on ISIC 2019/2020 or private clinical datasets would strengthen the generalizability claims.
    \item \textbf{Fixed image resolution:} All models used $224 \times 224$ input resolution. Higher resolutions (e.g., $384 \times 384$) may capture finer details relevant for subtle lesions but would increase computational cost.
    \item \textbf{No external clinical validation:} The models were not evaluated by practicing dermatologists in a clinical workflow, which would be necessary before any real-world deployment.
    \item \textbf{Limited Transformer exploration:} Only Swin-Tiny was evaluated. Larger Swin variants or other ViT architectures (e.g., DeiT, BEiT) might perform differently.
\end{itemize}

\section{Future Work}

Several promising directions for future research emerge from this study:

\begin{itemize}
    \item \textbf{Web-based deployment:} Develop an interactive web application (using Streamlit or Gradio) that allows clinicians to upload dermoscopy images and receive real-time predictions with Grad-CAM visualizations.
    \item \textbf{Test-Time Augmentation (TTA):} Apply multiple augmentations during inference and average the predictions, potentially boosting accuracy by 1--2\% without retraining.
    \item \textbf{Knowledge distillation:} Transfer the knowledge from the ensemble into a single lightweight model suitable for mobile or edge deployment.
    \item \textbf{Self-supervised pretraining:} Leverage large unlabeled dermoscopy image collections to pretrain models using contrastive learning (e.g., SimCLR, DINO) before fine-tuning on the labeled ISIC data.
    \item \textbf{Multi-task learning:} Jointly predict diagnosis and segmentation masks, as lesion boundary information may improve classification accuracy.
\end{itemize}
